\documentclass[a4paper]{article}
\usepackage[margin=3cm]{geometry}
\usepackage{graphicx}

\usepackage{hyperref}
\hypersetup{
	colorlinks=true,
	linkcolor=blue,
	filecolor=magenta,
	urlcolor=cyan,
	hyperindex=true,
	linktocpage=true,%makes the page number be the link in the index
}
\usepackage{listings}
\lstset{
	basicstyle=\tiny\ttfamily,
	columns=fullflexible,
	keepspaces=true,
	breaklines=true,
	frame=single,
	showstringspaces=true,
	tabsize=4,
	showspaces=true,
	showtabs=true,
}

\renewcommand{\arraystretch}{1.2}

\title{Felicity tech tree:\\Gallium Arsenide photovoltaics}

\author{Chaya2766}

\date{last updated \today}

\begin{document}
	\sffamily
	\hyphenchar\font=-1
	\sloppy
	
	\maketitle
	
	\textbf{
		\sloppy
		\hyphenchar\font=-1
		}
	
	\vfill
	
	\begin{abstract}
		Math behind the photovoltaic power generation in the Felicity tech tree.
		
		\medskip
		
		References an earlier \href{https://chaya2766.github.io/stareater-expanse-wiki/?page=documents/ion_beam_printing/ion_beam_printing.pdf}{"Ion beam printing" document by me}, as well as the paper \href{https://pmc.ncbi.nlm.nih.gov/articles/PMC8200097/pdf/materials-14-03075.pdf}{"Overview of the Current State of Gallium Arsenide-Based Solar Cells" by Nikola Papež, Rashid Dallaev, Stefan Talu and Jaroslav Kaštyl}.
	\end{abstract}
	
	%\noindent \rule{\linewidth}{1pt}
	\vfill
	
	\begin{quote}
		\begin{center}
			\textbf{DISCLAIMER AND WARNING}
		\end{center}
		
		This is not a scientific paper or engineering advice.
		This document is only intended as creative world‑building and has not been peer‑reviewed or approved for any practical application.
		Use of these ideas in real‑world designs or other serious application is entirely at the reader’s discretion and risk.
	\end{quote}
	
	\vfill
	
	\tableofcontents
	
	\vfill
	
	Reference markings:
	
	[1] - "Ion beam printing"
	
	[2] - "Overview of the Current State of Gallium Arsenide-Based Solar Cells"
	
	[3] - \href{https://en.wikipedia.org/wiki/Gallium_arsenide}{wikipedia page for "Gallium Arsenide"}
	
	[4] - "Felicity tech tree: low temperature electrolysis of mixed metal oxides"
	
	[5] - Engineering Toolbox table of emissivities for various materials \\ \indent (https://www.engineeringtoolbox.com/emissivity-coefficients-d\_447.html)
	
	\pagebreak
	
	\section{Characteristics of GaAs photovoltaics}
	
	GaAs photovoltaics are made 10$\mu$m thick [2], and achieve an efficiency of 40\%.
	
	\bigskip
	
	Given that Gallium Arsenide has a density of 5.3176 g/cm3 [3], the mass of material that goes into creating the panels would be 53.176g/m$^2$.
	
	$$ \frac{5.3176g/cm^3}{10\mu m} = 53.176g/m^2 $$
	
	Given that molar mass of Gallium Arsenide is 144.645 g/mol [3], the amount of material needed to create the panes is around 0.368mol of GaAs per square meter.
	This means 0.368mol of Gallium and 0.368mol of arsenic, as GaAs contains both in equal parts, giving a total of 0.735 mol/m$^2$ counting individual elements.
	For the individual elements:
	
	\begin{itemize}
		\item Gallium : 69.72 g/mol, results required 0.36763mol/m$^2$, 25.6312g/m$^2$
		
		\item Arsenic : 74.9216 g/mol, results required 0.36763mol/m$^2$, 27.5435g/m$^2$
	\end{itemize}
	
	
	$$ \frac{53.176g/m^2}{144.645 g/mol} = 0.3676310968 mol/m^2 \qquad 2 \cdot 0.3676310968 mol/m^2 = 0.7352621936 mol/m^2 $$
	
	$$ (0.3676310968 mol/m^2) \cdot (69.72 g/mol) = 25.63124007 g/m^2 (Ga) $$
	
	$$ (0.3676310968 mol/m^2) \cdot (74.9216 g/mol) = 27.54350998 g/m^2 (As) $$
	
	\bigskip
	
	The printer discussed in [3] is capable of depositing 5 micromol per second, and consumes 202.62MJ/mol.
	As such the panels would take around 147ks/m$^2$ and 149MJ/m$^2$ to print.
	
	$$ \frac{0.7352621936 mol/m^2}{5E-6 mol/s} = 147052.4387 s/m^2 $$
	
	$$ 0.7352621936 mol/m^2 \cdot 202.62MJ/mol = 148978825.7 J/m^2 $$
	
	\bigskip
	
	All taken together:
	
	\begin{center}
		\begin{tabular}{| c | c |}
			\hline
			Solar conversion efficiency & 40\% \\\hline
			mass & 25.6312g/m$^2$ Gallium \\
			& 27.5435g/m$^2$ Arsenic \\
			& total 53.176g/m$^2$ \\\hline
			material amount (single atoms) & 0.36763mol/m$^2$ Gallium,\\
			& 0.36763mol/m$^2$ Arsenic,\\
			& total 0.73526mol/m$^2$ \\\hline
			time to print & 147.052ks/m$^2$ \\\hline
			energy to print & 148.979MJ/m$^2$ \\\hline
		\end{tabular}
	\end{center}
	
	\pagebreak
	
	\section{Construction of concentrated solar arrays}
	
	GaAs photovoltaics can handle sunlight at intensities of 246 to 302 suns [2], as such, ignoring absorbtion for now, the area of the collector can be that many times smaller than the mirror area.
	
	\bigskip
	
	As a basic standard design I will assume a concentrated solar array where the collector is a disk 5m in radius, and the mirror as seen from the front is 50m in radius.
	
	\bigskip
	
	Covering a 5m-radius collector in GaAs photovoltaics would require printing a total area of 15.708m$^2$.
	
	This results in the following resource costs:
	
	\begin{itemize}
		\item Gallium: 5.775mol, 0.402615kg
		
		\item Arsenic: 5.775mol, 0.4326533kg
		
		\item Energy: 2340.162132 MJ
		
		\item Time: 2.340162132 Ms ($\approx$ 27 d + 2 h + 2 min + 42.132 s)
	\end{itemize}
	
	$$ M_{Ga} = 25.6312g/m^2 \cdot 15.708m^2 = 0.4026148896 kg (Ga) $$
	
	$$ N_{Ga} = 0.36763mol/m^2 \cdot 15.708m^2 = 5.77473204 mol (Ga) $$
	
	$$ M_{As} = 27.5435g/m^2 \cdot 15.708m^2 = 0.432653298 g (As) $$
	
	$$ N_{As} = 0.36763mol/m^2 \cdot 15.708m^2 = 5.77473204 mol (As) $$
	
	$$ E = 148.979MJ/m^2 \cdot 15.708m^2 = 2340.162132 MJ $$
	
	$$ T = 148.979ks/m^2 \cdot 15.708m^2 = 2340162.132 s$$
	
	The skeleton of the concentrated solar array has a volume of 78.7869m$^3$, measured empirically by rendering a 3d model (code at the end) and measuring the volume in Orca slicer.
	The assumed thickness for both the mirror, supports, and the collector was 1cm.
	
	\medskip
	
	Assuming metal density of 7750kg/m$^3$, the skeleton would mass 610598.475 kg.
	
	\medskip
	
	Given that the LTE process [4] can produce metal at a cost of 14.7625MJ/kg, and rate of 1.008kg/h, that amount of metal would take 9014GJ and 2.18071Gs to produce.
	
	(2.18071Gs $\approx$ 25239 d + 16 h + 27 min + 19 s)
	
	$$ 78.7869m^3 \cdot 7750kg/m^3 = 610598.475 kg $$
	
	$$ 610598.475 kg \cdot 14.7625MJ/kg = 9.013959987 TJ $$
	
	$$ \frac{610598.475 kg}{1.008kg/h} = 2180708839 s $$
	
	\pagebreak
	
	The concentrated solar array is collecting light with a mirror 50m in radius, but the collector which is 5m in radius is in front of the mirror, blocking some of the light.
	The mirror is also only 93\% reflective since smooth polished steel absorbs 7\% of light [5].
	The photovoltaics have an efficiency of 40\%.
	
	$$ (\pi (50m)^2 - \pi (5m)^2) \cdot 0.93 \cdot 0.4 = 7775.441818 m^2 \cdot 0.93 \cdot 0.4 = 2892.464356 m^2 $$
	
	The array is equivalent to a 100\% efficient collector with an area of $2892.464356 m^2$.
	
	\bigskip
	
	At earth's distance, sunlight intensity is 1361W/m$^2$, which means the array would produce 3.9366MW of power.
	
	$$ 2892.464356 m^2 \cdot 1361W/m^2 = 3936643.989 W $$
	
	This can be generalized for largely any level of insolation, for example, the asteroid belt, at the distance of 2.7AU, has an insolation of 186.69W/m$^2$, so the array would produce 540kW of power.
	
	$$ 1361W/m^2 \cdot (\frac{1AU}{2.7AU})^2 \cdot 2892.464356 m^2 = 186.6941015 W/m^2 \cdot 2892.464356 m^2 = 540006.0341 W $$
	
	And at Saturn's distance, 9.5826 AU, the insolation is only 14.8215W/m$^2$ so the array would produce 42.87kW of power.
	
	$$ 1361W/m^2 \cdot (\frac{1AU}{9.5826AU})^2 \cdot 2892.464356 m^2 = 14.82147429 W/m^2 \cdot 2892.464356 m^2 = 42870.58609 W $$
	
	\rule{\linewidth}{0.5pt}
	
	Cost of the array, all taken together:
	
	\begin{itemize}
		\item 610598.475 kg of mixed metal alloy in the mirror and support skeleton
		\item 0.402615kg (5.775mol) of Gallium in the photovoltaic cells
		\item 0.4326533kg (5.775mol) of Arsenic in the photovoltaic cells
		\item 2340.162132 MJ energy spent Ion beam printing the photovoltaic cells
		\item 2.340162132 Ms total Ion beam print time for the photovoltaic cells
	\end{itemize}
	
	Ofcourse additional energy must be spent producing the materials if they are not available:
	
	\begin{itemize}
		\item 9.013959987 TJ energy spent on electrolysis for the mixed metal alloy
		
		\item 2 180 708 839 s total electrolysis time for the mixed metal alloy
		
		\item todo: Germanium and Arsenic production costs, that will have to wait for a future document
	\end{itemize}
	
	\pagebreak
	
	\section{Concentrated solar array model}
	
	Here is the OpenSCAD code for the 3d model that was used to calculate the volume of the concentrated solar array.
	
	\begin{lstlisting}
module parabolic_mirror(focal_length=20,mirror_length=80,thickness=1,detail=16){
    // Parameters
    f = focal_length;           // focal length (mm)
    z_max = mirror_length;       // mirror depth (mm)
    segments = detail;   // increase for smoother curve
    thickness = thickness;    // rim thickness (optional)

    // build 2D profile points (x,z)
    pts = [for (i = [0:segments]) let(z = z_max * i/segments, x = 2*sqrt(f*z)) [x, z]];
    // optionally close inner vertex and add rim thickness
    profile = concat([[0,0]], pts, [[0,z_max+thickness]]);

    difference(){
        color("#888"){rotate_extrude($fn = segments)polygon(points = profile);}
        translate([0,0,thickness]){
            color("#CCC"){rotate_extrude($fn = segments)polygon(points = profile);}
        }
    }
}

//100,80 = 180
//50,80 = 130
//50,50 = 100
//so the radius of the mirror as seen from head on is just the focal length plus mirror length

module concentrated_solar_array(array_radius,collector_radius,thickness=1,detail=16){
    /*
    focus = array_radius*(3/4);
    depth = array_radius*(1.33/4);
    //1.33 derived empirically to achieve correct radius
    */
    focus = array_radius;
    depth = array_radius/4;
    //this new way of calculating the focus and depth also results in correct size
    //but this time it feels way too elegant to be a coincidence
    //I don't know the math behind why it should work out this way
    parabolic_mirror(focus,depth,thickness,detail);
    color("#888"){
    translate([0,0,focus]){cylinder(r=collector_radius,h=thickness);}
    for(i=[0:1:8]){
        rotate([0,0,i*45]){
            hull(){
                translate([array_radius-thickness,0,depth]){sphere(thickness/2);}
                translate([collector_radius,0,focus+thickness/2]){sphere(thickness/2);}
            }
        }
    }}
}

module mobile_concentrated_solar_array(array_radius,collector_radius,thickness=1){
    concentrated_solar_array(array_radius,collector_radius,thickness);
    translate([0,0,-array_radius/2]){cylinder(r=0.5,h=2*array_radius);}
    translate([0,0,-array_radius/2]){sphere(r=3);}
    translate([0,0,1.5*array_radius]){sphere(r=3);}
    }

rotate([0,90,0]){
concentrated_solar_array(50,5,0.01,32);
}

$vpr=[105,0,35];
$vpd=325;
	\end{lstlisting}
	
\end{document}